\chapter{对称性、简并性与角动量阶梯}
\providecommand{\ket}[1]{\left\lvert#1\right\rangle}

\section{对称性与简并能级}

对称性是一种不改变系统物理描述的运算。它不改变哈密顿量、Schr\"odinger方程或可能的能量。若用幺正算符\(S\)表示该运算，则不变性意味着
\begin{equation}
S H S^\dagger=H,
\end{equation}
或等价地
\begin{equation}
[S,H]=0.
\end{equation}

平移就是一个例子。把一个原子从一个地方移到另一个地方，其能量不变。这对每一个原子态都成立，而不仅是对基态。旋转是另一个例子，并且将是这里的主要主题。

交换两个全同粒子也是一种对称性。把两个电子称为电子一和电子二，或者哈里和弗雷德，并不会使它们变得不同。交换标签不能改变原子的性质。

晶体还有一种额外的平移对称性。把整个晶体在空间中移动是普通的平移。然而，对于无限晶体，平移一个晶格间距也会得到相同的物理排列。等价地，可以通过把每个格点指标移动一来重新标记格点。对于有限晶体，边界会破坏这一精确陈述，但在很长的晶体中远离两端处，它成为一个良好的对称性。

空间的均匀性---即每个位置都与其他任何位置相同---就是平移对称性。空间的各向同性---即每个方向都与其他任何方向相同---就是旋转对称性。

简并与对称性有关，但二者并不相同。当一个以上的态具有完全相同的能量时，该能级就是简并的。对于一个具有离散谱的一般哈密顿量，两个无关能级之间的精确相等将是惊人的巧合。符合到\(10^{15}\)分之一的精度仍然不是精确相等。精确简并通常指向对称性。

对称性有时蕴含简并，但并非总是如此。第一个例子既说明了为什么它可能如此，也说明了为什么单靠一种对称性可能不够。

\section{圆环上的粒子}

考虑一个被约束在圆形导线上的珠子。它的位置可以用笛卡尔坐标\(x\)和\(y\)来描述，但一个角坐标就足够了：
\begin{equation}
\psi=\psi(\theta).
\end{equation}
圆上的位置概率密度为\(\psi^*(\theta)\psi(\theta)\)。

一个逆时针方向的小角度\(\epsilon\)的主动旋转给出
\begin{equation}
\psi(\theta)\longrightarrow\psi(\theta-\epsilon).
\end{equation}
因此
\begin{align}
\delta\psi(\theta)
&=\psi(\theta-\epsilon)-\psi(\theta)\\
&=-\epsilon\frac{d\psi}{d\theta}
+\mathcal O(\epsilon^2).
\end{align}

算符
\begin{equation}
L_z=-i\hbar\frac{d}{d\theta}
\end{equation}
在周期波函数上是厄米的。它是
\begin{equation}
p_x=-i\hbar\frac{d}{dx}.
\end{equation}
的角类比。因子\(i\)对于厄米性是必不可少的。用\(L_z\)表示，
\begin{equation}
\delta\psi
=-\frac{i\epsilon}{\hbar}L_z\psi
+\mathcal O(\epsilon^2),
\end{equation}
所以无穷小旋转算符为
\begin{equation}
R(\epsilon)
=1-\frac{i\epsilon}{\hbar}L_z
+\mathcal O(\epsilon^2).
\end{equation}
这确认了\(L_z\)是绕垂直于圆环的轴的旋转生成元。

\begin{classroomqa}
\lecturetimestamp{00:11:19}
\audiencequestion{涉及\(i\)的因子之一是否应该取相反的符号？}
\lecturerresponse{检查符号得到上述关系。在约定\(\psi(\theta)\to\psi(\theta-\epsilon)\)下，变化为\(-\epsilon\,d\psi/d\theta\)，它等于\(-i\epsilon L_z\psi/\hbar\)。}
\end{classroomqa}

现在求角动量的本征态：
\begin{equation}
L_z\psi=M\psi.
\end{equation}
代入微分算符得到
\begin{equation}
-i\hbar\frac{d\psi}{d\theta}=M\psi,
\end{equation}
其解具有形式
\begin{equation}
\psi(\theta)=C\exp\left(\frac{iM\theta}{\hbar}\right).
\end{equation}

还有一个在无限直线上运动时不会出现的附加条件。点\(\theta=0\)和\(\theta=2\pi\)是同一点，所以单值轨道波函数必须满足
\begin{equation}
\psi(\theta+2\pi)=\psi(\theta).
\end{equation}
对于上述本征函数，这要求
\begin{equation}
\exp\left(\frac{2\pi iM}{\hbar}\right)=1.
\end{equation}
因此，
\begin{equation}
M=m\hbar,
\qquad
m\in\mathbb Z.
\end{equation}
所以圆环上可能的轨道角动量为
\begin{equation}
\ldots,-2\hbar,-\hbar,0,\hbar,2\hbar,\ldots.
\end{equation}

角动量为\(m\hbar\)的态是否必须与角动量为\(-m\hbar\)的态具有相同的能量？正负\(m\)描述相反的环流方向。问题在于，仅凭旋转对称性是否要求这两种运动具有相同的能量。

\begin{classroomqa}
\lecturetimestamp{00:22:51}
\audiencequestion{使用\(\psi(\theta-\epsilon)\)是否意味着粒子沿顺时针方向运动？}
\lecturerresponse{不。该约定只确定了哪个方向被称为正旋转。它定义了角位置和角动量的符号，但并不决定粒子实际运动的方向。}
\end{classroomqa}

一个垂直于圆环平面的磁场提供了一个反例。它可以使正负角动量的能量向相反方向移动，从而分裂这一对。移动的符号取决于电荷和取向约定，但分裂才是重点。

这样的场仍然保持绕垂直轴的旋转。因此它表明，仅凭旋转对称性并不要求
\begin{equation}
E(m)=E(-m).
\end{equation}
没有这种分裂时，\(m=1\)和\(m=-1\)可以构成一对简并态，\(m=2\)和\(m=-2\)也可以。态\(m=0\)没有这种意义上的不同伙伴。

\section{反射对称性与\texorpdfstring{\(m,-m\)}{m, -m}对}

加入一个离散反射对称性。令\(\mathcal M\)表示关于圆的一条直径的反射：
\begin{equation}
(\mathcal M\psi)(\theta)=\psi(-\theta).
\end{equation}
它反转环流的方向：
\begin{equation}
\mathcal M\psi_m=\psi_{-m}.
\end{equation}

若反射是一种对称性，则
\begin{equation}
[\mathcal M,H]=0.
\end{equation}
假设
\begin{equation}
H\ket{m}=E_m\ket{m}.
\end{equation}
施加\(\mathcal M\)得到
\begin{align}
H\mathcal M\ket{m}
&=\mathcal M H\ket{m}\\
&=E_m\mathcal M\ket{m}.
\end{align}
由于\(\mathcal M\ket{m}\)是角动量为\(-m\hbar\)的态，
\begin{equation}
E_{-m}=E_m.
\end{equation}
旋转对称性守恒角动量；反射对称性将两个相反的值联系起来。二者共同要求这一对态简并。

\begin{classroomqa}
\lecturetimestamp{00:29:54}
\audiencequestion{既然\(m\)和\(-m\)模具有相同的概率密度，这是否已经表明它们的能量相等？}
\lecturerresponse{不。波函数\(e^{im\theta}\)和\(e^{-im\theta}\)虽然绝对值平方相同，但本身不同。角动量本征值方程涉及波函数，而不仅仅是\(\psi^*\psi\)。量子力学需要振幅，而不仅仅是概率。}
\end{classroomqa}

磁场保持绕其轴的旋转，但破坏了上面所用的反射。

\begin{classroomqa}
\lecturetimestamp{00:33:05}
\audiencequestion{反射对称性会反转磁场吗？}
\lecturerresponse{会。磁场是轴矢量。反射一个指向板内的场会得到一个指向板外的场。等价地，一个电流回路及其镜像携带相反方向的电流，并产生方向相反的场。}
\end{classroomqa}

旋转与反射的非对易性可以直接看出。作用于角动量本征函数时，
\begin{align}
\mathcal M L_z\psi_m
&=m\hbar\,\psi_{-m},\\
L_z\mathcal M\psi_m
&=-m\hbar\,\psi_{-m}.
\end{align}
因此
\begin{equation}
[\mathcal M,L_z]\psi_m
=2m\hbar\,\psi_{-m},
\end{equation}
当\(m\neq0\)时它非零。先旋转再反射与先反射再旋转不是同一运算。

当两个变换都是对称性时，两个算符都与\(H\)对易，但它们彼此不对易。这正是导致简并的模式。

\begin{classroomqa}
\lecturetimestamp{00:40:28}
\audiencequestion{任何非对易的海森伯对，例如\(x\)和\(p\)，都会产生简并吗？}
\lecturerresponse{一般不。这些算符还必须是对称性：它们必须与哈密顿量对易。位置和动量通常并不都守恒。如果二者确实都与\(H\)对易，那么它们的非对易性将意味着简并。}
\end{classroomqa}

\section{对称性的对易子代数}

假设\(A\)和\(B\)是两个连续对称性的厄米生成元：
\begin{equation}
[A,H]=0,
\qquad
[B,H]=0.
\end{equation}
这些方程中没有任何东西要求\(A\)和\(B\)彼此对易。

对易子\([A,B]\)是反厄米的，因此定义一个厄米算符
\begin{equation}
C=i[A,B].
\end{equation}
数值因子和整体符号在这里无关紧要。要点是\(C\)也与哈密顿量对易。确实，
\begin{equation}
[AB,H]
=A[B,H]+[A,H]B=0,
\end{equation}
类似地
\begin{equation}
[BA,H]=0.
\end{equation}
因此
\begin{equation}
[C,H]
=i\bigl([AB,H]-[BA,H]\bigr)=0.
\end{equation}

新的生成元\(C\)可能独立于\(A\)和\(B\)，也可能只是已知生成元的线性组合。如果它是新的，就把它加入列表并构造进一步的对易子。重复这一过程可能无限继续，但往往在有限多个生成元之后闭合：每个进一步的对易子都是已有生成元的线性组合。这个闭合系统就是一个对易子代数或李代数。

有限的对称变换构成一个群。它们的无穷小生成元构成相应的李代数。

简并之所以重要，是因为它们组织谱。如果两个态能量相等，一个态就不能通过发射携带其能量差的光子衰变到另一个态。在原子光谱学中，能级的排列和简并决定可能的光谱线。同样的逻辑适用于激发态强子及其跃迁。

\begin{classroomqa}
\lecturetimestamp{00:49:07}
\audiencequestion{这些对易子本身是 symmetry group 的元素吗？}
\lecturerresponse{所讨论的算符是生成元。生成元及其对易子构成一个对易子代数，或李代数。由生成元构造的有限对称操作是群元素。}
\end{classroomqa}

如果所有对称性生成元彼此对易，则该对称性称为阿贝尔的。如果其中某些生成元的对易子不为零，则它是非阿贝尔的。

\begin{classroomqa}
\lecturetimestamp{00:50:25}
\audiencequestion{生成元的根本定义是什么？}
\lecturerresponse{取一个非常接近恒等的幺正对称变换，并将其写成\(1-i\epsilon G/\hbar+\cdots\)。厄米算符\(G\)就是它的生成元。重复相应的无穷小变换便可构建出有限变换。}
\end{classroomqa}

\begin{classroomqa}
\lecturetimestamp{00:52:08}
\audiencequestion{对称性生成元的代数是什么意思？}
\lecturerresponse{它是这样一组生成元的集合：它们可以相加、相减、乘以普通数，并可以取对易子。在这个代数中，对易子扮演乘法的角色，而封闭性意味着这些运算不会产生该集合之外的生成元。}
\end{classroomqa}

\section{绕不同轴的转动}

转动提供了非阿贝尔对称性最重要的例子。绕\(x\)轴的一次转动之后再绕\(y\)轴转动一次，一般并不等同于按相反顺序进行这两次转动。它们之间的差异本身与绕第三轴的转动有关。

\begin{classroomqa}
\lecturetimestamp{00:54:22}
\audiencequestion{既然有限转动不对易，那么无穷小转动对易吗？}
\lecturerresponse{它们相差在两个小转动角的二阶项上。这就足够了：把这两个小参数因子提出去之后，生成元本身具有非零的对易子。}
\end{classroomqa}

在计算量子对易子之前，先考虑相应的经典图像。取一个圆形的行星轨道，其角动量垂直于轨道平面。将整个轨道绕\(y\)轴转动，得到另一个允许的轨道，其能量完全相同。然而，它的角动量矢量现在具有不同的\(x\)分量。

因此，转动不变性产生了具有不同角动量分量的同能量位形。在量子力学中，这表现为绕不同轴的转动生成元之间的非对易性。

\section{笛卡儿坐标中的角动量}

首先考虑一个在\(xy\)平面内运动的粒子，而不把它限制在圆周上。在微小的逆时针旋转下，
\begin{equation}
\delta x=-\epsilon y,
\qquad
\delta y=\epsilon x.
\end{equation}
对于主动旋转后的态，变换后的波函数在逆旋转后的坐标处取值：
\begin{equation}
\psi_\epsilon(x,y)
=\psi(x+\epsilon y,y-\epsilon x).
\end{equation}
因此，
\begin{align}
\delta\psi
&=\epsilon\left(
y\frac{\partial}{\partial x}
-x\frac{\partial}{\partial y}
\right)\psi
+\mathcal O(\epsilon^2)\\
&=-\frac{i\epsilon}{\hbar}
\left(xp_y-yp_x\right)\psi
+\mathcal O(\epsilon^2).
\end{align}
于是，绕\(z\)轴的旋转的生成元为
\begin{equation}
L_z=xp_y-yp_x.
\end{equation}
这正是熟悉的经典矢量
\begin{equation}
\mathbf L=\mathbf r\times\mathbf p.
\end{equation}
的\(z\)分量。它不仅是力学角动量；在量子力学中，它是生成波函数旋转的算符。\footnote{选择主动旋转约定，是为了与先前的变换\(\psi(\theta)\mapsto\psi(\theta-\epsilon)\)以及口头推导中中间符号不确定性之后的最终识别\(L_z=xp_y-yp_x\)一致。}

循环置换\(x\to y\to z\to x\)便给出其他分量。\footnote{\(L_x\)和\(L_y\)的公式是通过把讲演中明确规定的循环替换\(x\to y\to z\to x\)应用于\(L_z=xp_y-yp_x\)而重构出来的。}
\begin{align}
L_x&=yp_z-zp_y,\\
L_y&=zp_x-xp_z,\\
L_z&=xp_y-yp_x.
\end{align}

对于转动不变的系统，这些算符与哈密顿量对易：
\begin{equation}
[H,L_x]=[H,L_y]=[H,L_z]=0.
\end{equation}
在中心力场中运动的粒子是标准例子。这些方程就是量子形式的角动量守恒。

\section{角动量李代数}

正则对易关系为
\begin{align}
[x_i,x_j]&=0,\\
[p_i,p_j]&=0,\\
[x_i,p_j]&=i\hbar\,\delta_{ij}.
\end{align}
坐标与属于不同方向的动量分量对易。因此，例如\(x\)与\(p_y\)对易，而\(x\)不与\(p_x\)对易。

现在计算
\begin{equation}
[L_x,L_y]
=
[yp_z-zp_y,\;zp_x-xp_z].
\end{equation}
大多数项都为零。两个非零贡献为
\begin{align}
[yp_z,zp_x]&=-i\hbar\,yp_x,\\
[zp_y,xp_z]&=i\hbar\,xp_y.
\end{align}
因此
\begin{align}
[L_x,L_y]
&=i\hbar\left(xp_y-yp_x\right)\\
&=i\hbar L_z.
\end{align}
循环置换坐标标号便给出
\begin{align}
[L_x,L_y]&=i\hbar L_z,\\
[L_y,L_z]&=i\hbar L_x,\\
[L_z,L_x]&=i\hbar L_y.
\end{align}

\begin{classroomqa}
\lecturetimestamp{01:14:50}
\audiencequestion{当对易子给出\(L_z\)时，我们是否仍然必须证明\(L_z\)不是\(L_x\)和\(L_y\)的线性组合？}
\lecturerresponse{是的。有义务确立这种独立性，而不是假定它。此处不进行直接的代数证明。}
\end{classroomqa}

所展示的对易子在\(L_x\)、\(L_y\)和\(L_z\)张成的空间上封闭：一旦把三者都包括进来，进一步的对易不会产生该空间之外的算符。这就是角动量李代数。连同
\begin{equation}
[H,L_i]=0,
\qquad i=x,y,z,
\end{equation}
这些关系式包含了关于能谱的大量信息，而不需要直接求解角向Schr\"odinger方程。

量子谐振子也采用类似的代数策略，其中产生和湮灭算符生成其能谱。这里，类似的算符将改变角动量的\(z\)分量，而不是能量。

选择\(z\)轴只是选择一种表示。没有什么使\(z\)在物理上比\(x\)或\(y\)更可取，但选择一个分量会给出方便的基。记
\begin{equation}
L_z\ket{m}=m\hbar\ket{m}.
\end{equation}
任务是确定\(m\)可以取哪些值，以及具有不同\(m\)的态之间如何联系。

\begin{classroomqa}
\lecturetimestamp{01:24:02}
\audiencequestion{当你把这个构造与谐振子比较时，你指的是量子谐振子吗？前一节课没有讲过这个。}
\lecturerresponse{是的，比较的对象是量子谐振子。目前的角动量论证是自足的，不需要谐振子的知识。这个类比以后也可以反过来读：谐振子使用类似的代数方法。}
\end{classroomqa}

\section{角动量的升与降}

定义
\begin{equation}
L_+=L_x+iL_y,
\qquad
L_-=L_x-iL_y.
\end{equation}
角动量对易子给出
\begin{equation}
[L_z,L_+]=\hbar L_+,
\qquad
[L_z,L_-]=-\hbar L_-.
\end{equation}

把第一个关系式作用于\(L_z\)的一个本征态：
\begin{align}
L_zL_+\ket{m}
&=\left(L_+L_z+[L_z,L_+]\right)\ket{m}\\
&=(m+1)\hbar\,L_+\ket{m}.
\end{align}
除非\(L_+\ket{m}=0\)，它是\(L_z\)的另一个本征态，现在的本征值为\((m+1)\hbar\)。类似地，
\begin{equation}
L_zL_-\ket{m}
=(m-1)\hbar\,L_-\ket{m}.
\end{equation}
因此，\(L_+\)把\(m\)升高一，而\(L_-\)把它降低一。

从任意一个本征值出发，代数生成一个以整数相隔的序列。该序列可以无限延续，也可以终止，因为在上端
\begin{equation}
L_+\ket{m_{\mathrm{top}}}=0
\end{equation}
在下端
\begin{equation}
L_-\ket{m_{\mathrm{bottom}}}=0
\end{equation}

\begin{classroomqa}
\lecturetimestamp{01:35:11}
\audiencequestion{圆周论证不是已经要求角动量为整数吗？}
\lecturerresponse{那个论证证明了由单值位置波函数表示的轨道角动量取整数值。抽象的角动量代数给出一个更宽松的结果：阶梯既可以包含整数值，也可以包含半整数值。}
\end{classroomqa}

如果阶梯终止，旋转不变性把它的两端联系起来。绕垂直于\(z\)的轴旋转\(180^\circ\)，使
\begin{equation}
L_z\longrightarrow-L_z.
\end{equation}
因此，谱必须关于零对称：
\begin{equation}
m_{\mathrm{bottom}}=-m_{\mathrm{top}}.
\end{equation}
由于相邻值相差一，有两种可能。阶梯要么包含整数，
\begin{equation}
\ldots,-2,-1,0,1,2,\ldots,
\end{equation}
要么包含半整数，
\begin{equation}
\ldots,-\frac{5}{2},-\frac{3}{2},-\frac{1}{2},
\frac{1}{2},\frac{3}{2},\frac{5}{2},\ldots,
\end{equation}
端点位于升算符或降算符给出零的地方。

圆周上单值波函数的论证为轨道角动量选择了整数可能性，
\begin{equation}
\mathbf L=\mathbf r\times\mathbf p.
\end{equation}
半整数可能性也是物理的，属于自旋角动量，将另行处理。

最后，当哈密顿量具有旋转不变性时，阶梯态具有相同的能量。假设
\begin{equation}
H\ket{m}=E\ket{m}.
\end{equation}
由于每个分量\(L_i\)都与\(H\)对易，\(L_+\)和\(L_-\)也同样如此。因此
\begin{align}
H L_+\ket{m}
&=L_+H\ket{m}\\
&=E L_+\ket{m}.
\end{align}
只要\(L_+\ket{m}\neq0\)，升高后的态就与原态具有相同的能量。同样的计算给出
\begin{equation}
H L_-\ket{m}=E L_-\ket{m}.
\end{equation}
因此，升降操作生成了一族\(L_z\)本征值不同、能量却完全相同的态。这就是本讲始终在阐明的非对易对称性与简并之间的联系。

\begin{classroomqa}
\lecturetimestamp{01:44:13}
\audiencequestion{这些阶梯态是否与选择不同的旋转轴有关？}
\lecturerresponse{它们通过旋转相互联系。\(L_x\)或\(L_y\)的本征矢量是所选\(L_z\)本征矢量的线性组合。旋转一个态会改变其角动量分量，同时保持其能量不变。}
\end{classroomqa}

选择\(L_+\)和\(L_-\)是方便的，因为这些组合直接暴露出阶梯结构。这个技巧在量子力学中反复出现：当一个对易子代数闭合时，升算符和降算符往往能组织其本征值谱。

\begin{classroomqa}
\lecturetimestamp{01:46:33}
\audiencequestion{当\(L_+\)作用于最高态时会发生什么，顶端又是如何定义的？}
\lecturerresponse{它给出零。这就是终止阶梯顶端的定义。同样，\(L_-\)在底端给出零。抽象的阶梯不必终止。对于后面考虑的原子角动量态，人们通常会遇到终止的阶梯。}
\end{classroomqa}

下一步是进一步发展这种角动量结构，并将其应用于原子，包括氢原子。这里发现的旋转多重态解释了简并模式的一部分；进一步的结构仍有待展开。
